\scanpage{76}

% top, introduction
Posant

% top, DR and its graded family
\[
\unsure{\mathrm{DR}|\mathbb D\mathbb R}^{pq}=\mathrm{DR}^{pq}(X_*)=\operatorname{Ker}(C^{p,q+1}\xrightarrow{\bar T_{\unsure{\wedge|\bullet}}}C^{p,q+2})\ ,\qquad
\mathrm{DR}^{**}=(\mathrm{DR}^{pq})_{p,q\in\unsure{\mathbb Z|\mathbb Z\times\mathbb Z}}
\]

% upper, transition
on aura donc

% upper, two stages of the map and inclusions
\[
\underset{\text{épi si }p\ne0}{\smash[b]{\mathstrut C^{p,q}\xrightarrow{\bar T_.\,=u}\underset{\substack{\cap\\C^{p,q+1}}}{\mathrm{DR}^{pq}}}}\qquad
\text{d'où}\quad C^{**}\xrightarrow{u}\underset{\substack{\cap\,i\\ C_*^{**}}}{\mathrm{DR}^{**}}\quad\deleted{\unsure{u}}
\]

% upper left, boxed marginal identities
\margin{$\begin{aligned}\operatorname{Im}i&=\operatorname{Ker}u\\i\circ u&=\bar T_\wedge(\,.\,)\end{aligned}$}

% upper, (C)
\[
i(x).i(y)=0\tag*{(C)}
\]

% upper, parenthesis and degree restriction
(il est évident que $i(\mathrm{DR})\underline{i(u\,\mathrm{DR})=0}$ donc $i(\mathrm{DR}^{p,q})\,i(\mathrm{DR}^{p',q'})=0$ si $p$ ou $p'\ne0$)

% left, oblique boxed marginal reformulation
% slanted
\margin{donc (C) équiv$^{\mathrm t}$ à $\boxed{i(\mathrm{DR}^{0p}).i\mathrm{DR}_*^{0q}\mathrel{\unsure{=|\ne}}0}$}

% upper middle, algebra structure
\deleted{et} on aura sur $\mathrm{DR}^{**}$ une structure d'algèbre graduée différentielle à p.d. \inserted{canonique, avec axiomes usuels (alternée etc)} \deleted{\unsure{telle que}}

% upper middle, map u
\[
u:C^{**}\mathrel{\unsure{\relbar|\longrightarrow}}\unsure{D|\mathrm{DR}}^{**}
\]

% middle, compatibility and qualifications
soit compatible avec ttes ces structures \inserted{et $\left\{\begin{aligned}i(xuy)&=i(x)y\\i(uyx)&=y\,i(x)\end{aligned}\right.$}. Ces structures sur $\mathrm{DR}^{**}$ sont presque déterminées par ces propriétés, plus celle que l'inclusion $\mathrm{DR}^{**}\subset C^{**}$ est de bidegré $0,1$ et \inserted{anti}commute à la différentielle -- à la détermination près des produits $xy$ \deleted{\unsure{et $x$}} et des puissances divisées $x^{(i)}$ $(i\geq2)$ quand $x,y$ \deleted{sont} de degré complémentaire $0$, (auquel cas on ne peut pas néc. les écrire sous la forme $u(x')$\inserted{\deleted{\unsure{et}} ou $u(y')$}, \unsure{où|tiré|tels} dans $C^{**}$).

% middle lower, bracketed NB and crowded revisions
[ \uline{NB} Il est évident que sur \deleted{\unsure{$\mathrm{DR}$}} on peut définir \inserted{de façon unique} une multiplication \inserted{bilinéaire sur $\mathrm{DR}\times\mathrm{DR}$,\inserted{$(\operatorname{Im}u)\times\mathrm{DR}$ et $\mathrm{DR}\times\operatorname{Im}u$}\deleted{\unsure{$\in k$}}} \deleted{\unsure{et}} satisfaisant

% lower middle, three identities
\[
u(x.y)=u(x)u(y),\qquad i(xu(y))=i(x)y,\quad i(u(y)x)=y\,i(x)
\]

% lower middle, implication and applying i
(\deleted{\unsure{car}} la première relation équivaut, en appliquant $i$, à $\bar T_.(xy)=i(ux.uy)$, i.e. en appliq. la deuxième $\bar T.(xy)=(\bar T.x)y$

% lower, two maps and existence
donc il suffit de vérifier l'existence et l'unicité de $\operatorname{Im}u\times\mathrm{DR}\relbar\mathrm{DR}$ \inserted{et $\mathrm{DR}\times\operatorname{Im}u\relbar\mathrm{DR}$} satisfaisant les rel. 2 et 3. Or pour la $2^{\mathrm{ème}}$, il suffit de noter que

% lower, a)
a) $i(x)y$\inserted{$=xy$} $\in\deleted{\operatorname{Ker}\bar T}\mathrm{DR}=\operatorname{Ker}(\bar T_\wedge)$ car $\bar T(xy)=(\bar T x)y=0$ car $\bar T x=0$ puisque $x\in\mathrm{DR}$

% foot, b)
b) $i(x)y$\inserted{$=xy$} ne dépend que de $x\in\mathrm{DR}$ et de $uy\in\mathrm{DR}$, car si $\deleted{\unsure{u}}u\unsure{y)|(y)}=0$ i.e. $\bar Ty=0$, i.e. $y\in\mathrm{DR}$, on a $xy=0$ -- Idem pour 3.

% lower left
% archive sheet number: 75
